\documentclass[a4paper]{article}

\usepackage{graphicx}
\usepackage{amsmath,amssymb}
\usepackage{minted}
\usepackage{multirow}
\usepackage{float}
\usepackage{todonotes}
\usepackage{forest}
\usepackage{xstring}
\usepackage{xcolor}
\usepackage[most]{tcolorbox}
\usepackage{xparse}
\usepackage{natbib}
\usepackage{tikz}

\usetikzlibrary{graphs,graphdrawing}
\usegdlibrary{layered}

% for external graphviz
\input{/home/donkarlo/Dropbox/repo/nd_ai_project/src/nd_ai/knoweldge_representation/language/formal/computer/programming/tex/latex/code_clips/external_graphviz.tex}

% for tags
\input{/home/donkarlo/Dropbox/repo/nd_ai_project/src/nd_ai/knoweldge_representation/language/formal/computer/programming/tex/latex/parts/control_sequence/kind/semtag/semtag.tex}

% for links
\input{/home/donkarlo/Dropbox/repo/nd_ai_project/src/nd_ai/knoweldge_representation/language/formal/computer/programming/tex/latex/code_clips/seehere.tex}

\title{Particle Filtering in Discrete State-Space Models}
\date{}

\begin{document}
    \maketitle
    
    \cite{chopin-2001-sequential-inference-and-state-number-determination-for-discrete-state-space-models-through-particle-filtering} applies particle filtering to discrete state-space models, especially hidden Markov models with unobserved regime/state sequences. The paper proposes a Monte Carlo HMM filter that marginalizes hidden states to reduce particle-weight volatility and improve sequential Bayesian inference. It also introduces sequential state-number determination, estimating how many distinct hidden regimes have appeared over time.
    
    
    \section{Discrete State-Space Model}
        
        In a discrete state-space model, the observed signal \(y_t\) is assumed to be generated by an unobserved discrete state \(z_t\). If the hidden state at time \(t\) is \(k\), the observation model can be written as
        
        \[
            y_t
            \mid
            z_t = k
            \sim
            f_{\xi_k}(y_t),
        \]
        
        where \(f_{\xi_k}\) is the observation distribution associated with discrete state \(k\).
        
        The hidden state evolves in a discrete state space according to
        
        \[
            \Pr
            \left(
                z_t = l
                \mid
                z_{t-1} = k
            \right)
            =
            p_{kl},
        \]
        
        where \(p_{kl}\) is the transition probability from state \(k\) to state \(l\).
    
    
    \section{Particle Approximation}
        
        Particle filtering approximates the posterior distribution over the unknown parameters and hidden states by a set of weighted particles:
        
        \[
            \pi_t
            =
            \pi
            \left(
                \theta,
            z_{1:t}
                \mid
                y_{1:t}
            \right)
            \approx
            \left\{
                \left(
                    \theta^{(j)},
                z_{1:t}^{(j)},
                w_t^{(j)}
                \right)
            \right\}_{j=1}^{H}.
        \]
        
        Here, \(\theta\) denotes the unknown model parameters, \(z_{1:t}\) is the hidden state sequence, \(y_{1:t}\) is the observed sequence, \(w_t^{(j)}\) is the weight of particle \(j\), and \(H\) is the number of particles.
        
        When a new observation \(y_t\) arrives, each particle is updated according to how well it explains the new observation:
        
        \[
            w_t^{(j)}
            \propto
            w_{t-1}^{(j)}
            p
            \left(
                y_t
                \mid
                y_{1:t-1},
                \theta^{(j)}
            \right).
        \]
    
    
    \section{Marginalization of Discrete Hidden States}
        
        In the Monte Carlo HMM filter, the discrete hidden states are marginalized out. Therefore, the likelihood of the new observation is computed by summing over all possible discrete states:
        
        \[
            p
            \left(
                y_t
                \mid
                y_{1:t-1},
                \theta^{(j)}
            \right)
            =
            \sum_{k=1}^{K}
            \Pr
            \left(
                z_t = k
                \mid
                y_{1:t-1},
                \theta^{(j)}
            \right)
            f_{\xi_k^{(j)}}(y_t).
        \]
        
        This means that each particle does not need to carry one sampled hidden-state trajectory. Instead, it carries a probability distribution over the discrete states.
        
        The filtering probability of each state can then be estimated by the weighted average
        
        \[
            \widehat{\Pr}
            \left(
                z_t = k
                \mid
                y_{1:t}
            \right)
            =
            \frac{
                \sum_{j=1}^{H}
                w_t^{(j)}
                \Pr
                \left(
                    z_t = k
                    \mid
                    y_{1:t},
                    \theta^{(j)}
                \right)
            }{
                \sum_{j=1}^{H}
                w_t^{(j)}
            }.
        \]
    
    
    \section{General Interpretation}
        
        The method combines three elements:
        
        \[
            \text{state transition probabilities}
            \quad
            p_{kl},
        \]
        
        \[
            \text{observation likelihoods}
            \quad
            f_{\xi_k}(y_t),
        \]
        
        and
        
        \[
            \text{particle weights}
            \quad
            w_t^{(j)}.
        \]
        
        Thus, particle filtering in a discrete state-space model estimates the current hidden regime by combining how states evolve, how well each state explains the observation, and how probable each parameter particle is.
        
        \bibliographystyle{plainnat}
        \bibliography{/home/donkarlo/Dropbox/repo/nd_shared_research/bibliography/bibliography.bib}
\end{document}